% The boundary of an open annulus should be shown as dashed circles, not solid circles. This indicates that the boundary is not contained in the set being drawn.
% A puncture is usually marked as a hollow circle not a solid circle.
% The angular directions should really be between 0 and 360 degrees, not just between 0 and 180.
% It is better for a visual representation of the annulus that only the parts of the radial lines which are within the annulus are shown.
% I commented out the sample points which, in my opinion, created visual clutter.
% Instead of the sample points I added color-coded concentric circles around the origin in the punctured plane.
% The images of the concentric circles are shown in the annulus - the colors match.
% The correct notation for the new-style arrow tips in TikZ is like this: \draw[-{Latex}, ... ]
%
%
%\documentclass{standalone}
%\usepackage{tikz}
%\usetikzlibrary{decorations.pathreplacing, arrows.meta}
%\usepackage{amsfonts}
%
%\begin{document}
%
\begin{tikzpicture}[declare function = {F(\r) = (\r+2)/(\r+1);}, 
	baseline={(0,0)}]

% Left diagram: Punctured Euclidean plane
\begin{scope}[shift={(-6,0)}]  % Shift left by 6 units
    % Draw concentric circles to represent the punctured plane
    \draw[thin,orange] (0,0) circle (0.5);
    \draw[thin,green] (0,0) circle (1);
    \draw[thin,cyan] (0,0) circle (1.5);
    \draw[thin,blue] (0,0) circle (2);
    \draw[thin,violet] (0,0) circle (2.5);
    
    % Draw radial lines to represent angular directions
    \foreach \angle in {0, 30, 60, 90, 120, 150, 180, 210, 240, 270, 300, 330} {
	    \draw[gray, thin] (\angle:2pt) -- (\angle:3);
    }
    
    % Label the plane and mark the origin as a puncture
    \node[above] at (-1.5, 2.75) {$\mathbb{R}^2 \setminus \{(0,0)\}$};
    \draw[thick, red] (0,0) circle (2pt) node[below right] {$O$};  % Origin

%    % Sample points
%    \filldraw[blue] (0.5,0) circle (2pt) node[below left] {$(0.5, 0)$};
%    \filldraw[blue] (1,0) circle (2pt) node[below left] {$(1, 0)$};
%    \filldraw[blue] (2.5,0) circle (2pt) node[below left] {$(2.5, 0)$};

\end{scope}

% Right diagram: Open annulus
\begin{scope}[shift={(3,0)}]  % Shift right by 3 units
    % Draw the annulus
    \draw[thick, dashed] (0,0) circle (1);    % Inner circle
    \draw[thick, dashed, red] (0,0) circle (2);    % Outer circle
    
    % Draw the images of the concentric circles in the punctured plane, using the color for visual matching
	\draw[thin,orange] (0,0) circle ({F(0.5)});
	\draw[thin,green] (0,0) circle ({F(1)});
	\draw[thin,cyan] (0,0) circle ({F(1.5)});
	\draw[thin,blue] (0,0) circle ({F(2)});
	\draw[thin,violet] (0,0) circle ({F(2.5)});

    % Draw radial lines to show angular directions
    \foreach \angle in {0, 30, 60, 90, 120, 150, 180, 210, 240, 270, 300, 330} {
        \draw[gray, thin] (\angle:1) -- (\angle:2);
    }

    % Label the annulus
    \node[above] at (-1, 2.3) {$\{ (x, y): 1 < \sqrt{x^2 + y^2} < 2 \}$};
    
%    % Mark images of sample points
%    \filldraw[blue] (2,0) circle (2pt) node[below right] {$(2, 0)$};  % Image of (0.5, 0)
%    \filldraw[blue] (1.5,0) circle (2pt) node[below right] {$(1.5, 0)$};  % Image of (1, 0)
%    \filldraw[blue] (1.1,0) circle (2pt) node[below right] {$(1.1, 0)$};  % Image of (2.5, 0)
\end{scope}

% Draw an arrow representing the mapping
\draw[-{Latex}, thick, red] (-2.7, 0) -- (0.7, 0) node[midway, above] {$f(r, \theta) = \left( \frac{r+2}{r+1}, \theta \right)$};

% % Draw mapping for each sample point
% \draw[-{Latex}, thick, blue] (-5.5, 0) -- (3.5, 0);  % Mapping for (0.5, 0) -> (2, 0)
% \draw[-{Latex}, thick, blue ] (-5, 0) -- (3, 0);      % Mapping for (1, 0) -> (1.5, 0)
% \draw[-{Latex}, thick, blue] (-3.5, 0) -- (1.6, 0);  % Mapping for (2.5, 0) -> (1.1, 0)

\end{tikzpicture}%
%
%\end{document}